\documentclass[11pt,a4paper]{article}

\usepackage[T1]{fontenc}
\usepackage[utf8]{inputenc}
%\usepackage{lmodern}
\usepackage[margin=25mm]{geometry}
\usepackage{amsmath}
\usepackage{booktabs}
\usepackage{siunitx}
%\usepackage{microtype}
\usepackage[hidelinks]{hyperref}

\sisetup{per-mode=symbol, range-phrase={\text{--}}, range-units=single}

\title{Snubbers in Paralleled IGBT Modules\\
\large Application Notes for the ABB/Hitachi 5SNA~1000G650300 (\SI{6.5}{kV} / \SI{1000}{A} HiPak)}
\author{}
\date{\today}

\begin{document}
\maketitle

\begin{abstract}
These notes discuss the role of snubber networks when high-voltage IGBT modules are
operated in parallel, with reference to the ABB/Hitachi 5SNA~1000G650300 HiPak module
(\SI{6.5}{kV}, \SI{1000}{A}, SPT\textsuperscript{++} chip-set). The main conclusion is
that snubbers are useful in a paralleled configuration, but for a specific reason: they
act primarily as \emph{local decoupling elements} that shorten and equalize the
commutation loops of the individual modules, thereby limiting turn-off overvoltage and
improving dynamic current sharing. They do not correct static current imbalance, which
must be addressed by device matching, gate-drive design, and layout. Sizing guidelines
and a worked numerical example are given.
\end{abstract}

\section{Scope}
Paralleling of high-power IGBT modules is common practice when the required phase
current exceeds the rating of a single module. This document addresses one specific
design question: whether, and how, snubber networks should be used in a paralleled
configuration. The quantitative data refer to the 5SNA~1000G650300~\cite{ds}, but the
methodology applies to any high-voltage HiPak-class module. The general sizing
procedure (voltage budget, loop energy) follows the approach of the companion snubber
design guideline~\cite{snubnote}.

\section{Device Overview}
Table~\ref{tab:device} summarizes the datasheet values relevant to the present
discussion~\cite{ds}. Two figures deserve attention. First, the turn-off is
intentionally slow ($R_{G,\mathrm{off}}=\SI{15}{\ohm}$, $t_f \approx
\SIrange{380}{500}{ns}$), which limits the turn-off $di/dt$ of the IGBT itself. Second,
the anti-parallel diode exhibits a large reverse recovery ($I_{rr}$ up to
\SI{2.49}{kA}, $Q_{rr}$ up to \SI{2.26}{mC} at \SI{150}{\celsius}); at low temperature
and low forward current the recovery is snappier, and this is normally the most
violent transient in the commutation loop.

\begin{table}[htb]
\centering
\caption{Selected 5SNA~1000G650300 data (typical unless noted)~\cite{ds}.}
\label{tab:device}
\begin{tabular}{llll}
\toprule
Parameter & Symbol & Value & Conditions \\
\midrule
Blocking voltage & $V_{CES}$ & \SI{6500}{V} & $T_{vj} \geq \SI{25}{\celsius}$ \\
DC collector current & $I_C$ & \SI{1000}{A} & $T_C = \SI{110}{\celsius}$ \\
Fall time & $t_f$ & \SIrange{380}{500}{ns} & \SI{3600}{V}, \SI{1000}{A} \\
Turn-off delay & $t_{d(\mathrm{off})}$ & \SIrange{5.0}{5.9}{\micro s} & $R_G = \SI{15}{\ohm}$ \\
Diode peak recovery current & $I_{rr}$ & \SIrange{1.71}{2.49}{kA} & \SI{25}{}--\SI{150}{\celsius} \\
Recovered charge & $Q_{rr}$ & \SIrange{1.21}{2.26}{mC} & \SI{25}{}--\SI{150}{\celsius} \\
Short-circuit current & $I_{SC}$ & \SI{4800}{A} & $V_{CC} = \SI{4500}{V}$, $t_{psc} \leq \SI{10}{\micro s}$ \\
Internal stray inductance & $L_{\sigma,\mathrm{int}}$ & \SI{18}{nH} & CE terminals \\
RBSOA limit & $V_{CEM,\mathrm{chip}}$ & $\leq \SI{6500}{V}$ & at chip level \\
External gate--emitter capacitor & $C_{GE}$ & \SI{220}{nF} & datasheet test condition \\
\bottomrule
\end{tabular}
\end{table}

A package-level remark: the module exposes \emph{three} collector--emitter terminal
pairs that must be paralleled by the busbar. The module is therefore already an
internal parallel connection closed externally, and the busbar symmetry discipline
required between modules is equally required across the terminal pairs of a single
module.

\section{The Commutation Loop with Paralleled Modules}
With $N$ modules in parallel, the commutation loop splits into a \emph{common} section
(DC-link capacitor bank, shared laminated busbar) and $N$ \emph{branch} sections
(module terminals, local busbar tabs). Two effects make the paralleled case harder
than the single-module case:

\begin{enumerate}
\item \textbf{Summed $di/dt$ in the common section.} Each module turning off
\SI{1000}{A} in $t_f \approx \SI{0.5}{\micro s}$ contributes roughly
\SIrange{2}{2.5}{kA\per\micro s}. The common busbar section therefore sees
$N$ times this slope: \SIrange{4}{5}{kA\per\micro s} for two modules. The inductive
overvoltage of the common section scales accordingly.
\item \textbf{Asymmetry.} Any inequality between branch inductances produces unequal
transient overvoltages and unequal dynamic current sharing during turn-off and,
especially, during diode reverse recovery. The module with the lower-inductance branch
takes more of the commutated current and sees a different voltage trajectory.
\end{enumerate}

\section{What the Snubber Does --- and Does Not Do}
\subsection{Decoupling function}
A low-inductance film capacitor mounted directly across the DC terminals of
\emph{each} module (or each half-bridge leg, for side-by-side legs) makes the
effective commutation loop of every module short and, if the mounting is identical,
equal for all modules. The consequences are:
\begin{itemize}
\item the summed $di/dt$ is confined to the short local loop, and the common busbar
inductance is largely removed from the switching transient;
\item all modules commutate against the same local impedance, which improves dynamic
current sharing at turn-off and during diode recovery;
\item the turn-off overvoltage at the chip is reduced.
\end{itemize}
This decoupling function is the primary reason to use snubbers in a paralleled
configuration, and it is the arrangement recommended by the manufacturer's paralleling
application note~\cite{an2098}.

\subsection{Clamping function}
If a hard limit on the peak voltage is required --- in particular for short-circuit
turn-off, where the interrupted current is $I_{SC} = \SI{4800}{A}$ per module and the
loop overvoltage scales accordingly --- an RCD clamp per module is preferred over a
plain capacitor. The clamp absorbs the loop energy only above the DC-link voltage and
does not participate in the normal resonant exchange.

\subsection{What snubbers do not solve}
Snubbers act on the power loop only. They do not correct:
\begin{itemize}
\item \textbf{Static current sharing.} This is governed by $V_{CE,\mathrm{sat}}$
matching (same production batch) and thermal coupling (common heatsink). The
SPT\textsuperscript{++} chip-set has a positive temperature coefficient of
$V_{CE,\mathrm{sat}}$ at rated current, which is inherently self-balancing.
\item \textbf{Gate-side behavior.} Each module requires an individual gate resistor
on its own gate lead --- a shared gate resistor invites parasitic oscillations between
the paralleled gates --- plus the external $C_{GE} = \SI{220}{nF}$ per module as in the
datasheet test conditions, and a symmetric gate-loop layout with matched propagation
delays.
\end{itemize}

\section{Chip-Level Voltage Budget}
The RBSOA limit is specified at chip level. The chip voltage exceeds the terminal
voltage by the drop across the internal stray inductance:
\begin{equation}
V_{CE,\mathrm{chip}} = V_{CE,\mathrm{term}} + L_{\sigma,\mathrm{int}}\,\frac{di}{dt},
\qquad L_{\sigma,\mathrm{int}} = \SI{18}{nH}.
\end{equation}
The total peak chip voltage during turn-off is
\begin{equation}
V_{CE,\mathrm{chip}} = V_{DC}
+ \left(L_{\mathrm{common}}\,N + L_{\mathrm{branch}} + L_{\sigma,\mathrm{int}}\right)
\frac{di}{dt}\bigg|_{\mathrm{module}}
+ \Delta V_{\mathrm{snap}},
\label{eq:budget}
\end{equation}
where $\Delta V_{\mathrm{snap}}$ accounts for the diode snap-off contribution and
$L_{\mathrm{common}}$ is referred to a single module (its $di/dt$ is multiplied by $N$
because it carries the summed current).

One caution on the available margin: \SI{6.5}{kV} modules are typically operated from
a \SIrange{3.6}{4}{kV} DC link primarily because of cosmic-ray failure-rate
derating~\cite{an2042}, not to leave room for overshoot. The margin between $V_{DC}$
and $V_{CEM,\mathrm{chip}}$ is therefore not entirely available as an overshoot
allowance; the design budget should be set deliberately, typically reserving the
larger share for the short-circuit turn-off case.

\section{Sizing Guidelines}
\subsection{Decoupling capacitor}
The capacitor must absorb the magnetic energy of the loop section it decouples with an
acceptable voltage excursion $\Delta V$. From
$\tfrac{1}{2} L I^2 = \tfrac{1}{2} C \,\Delta V^2$:
\begin{equation}
C \;\geq\; L_{\mathrm{loop}} \left(\frac{I_{\mathrm{off}}}{\Delta V}\right)^{2}.
\label{eq:csize}
\end{equation}
The capacitor must be a low-inductance film type (dry PP), mounted directly on the
module terminals with identical hardware on every module; its own ESL adds to
$L_{\mathrm{branch}}$ and must be included in Eq.~\eqref{eq:budget}.

\subsection{Worked example}
Two modules, $V_{DC} = \SI{3600}{V}$, per-module turn-off current \SI{1000}{A},
$t_f = \SI{0.5}{\micro s}$ (hence $di/dt \approx \SI{2}{kA\per\micro s}$ per module),
$L_{\mathrm{common}} = \SI{100}{nH}$, $L_{\mathrm{branch}} = \SI{50}{nH}$.

Without snubbers, Eq.~\eqref{eq:budget} gives
\begin{equation}
\Delta V = (2 \cdot 100 + 50 + 18)\,\si{nH} \times \SI{2}{kA\per\micro s}
\approx \SI{536}{V},
\end{equation}
i.e.\ a peak chip voltage of about \SI{4.14}{kV} --- acceptable in nominal operation,
but the same loop interrupting $I_{SC} = \SI{4800}{A}$ per module with a faster
current collapse leaves considerably less margin, and the diode snap adds on top.

With a decoupling capacitor per module, the common \SI{100}{nH} is removed from the
transient. Allocating $\Delta V = \SI{300}{V}$ to the decoupled section at
$I_{\mathrm{off}} = \SI{1000}{A}$, Eq.~\eqref{eq:csize} yields
\begin{equation}
C \;\geq\; \SI{100}{nH} \times \left(\frac{\SI{1000}{A}}{\SI{300}{V}}\right)^{2}
\approx \SI{1.1}{\micro F}.
\end{equation}
Capacitors in the \SIrange{1}{4}{\micro F} class per module are typical for this
module family; a larger value relaxes $\Delta V$ at the cost of heavier inter-module
resonances (next section).

\subsection{RCD clamp}
If clamping is required, the clamp capacitor is sized with Eq.~\eqref{eq:csize} using
the short-circuit current and the clamp voltage rise as $\Delta V$; the discharge
resistor is chosen so that $\tau = R\,C$ is a small fraction of the switching period,
and its power rating from $P = \tfrac{1}{2} C \,\Delta V^2 f_{sw}$ per clamping event
rate. The clamp diode must be a fast soft-recovery type rated for the full DC-link
voltage.

\section{Inter-Module Resonances}
Plain capacitors on paralleled modules create LC meshes from module to module: the
snubber capacitors of adjacent modules resonate through the busbar inductance between
them. If the capacitors, their ESL, or their mounting are not identical, the result is
ringing and circulating currents flowing through the snubbers after every commutation.
Countermeasures, in order of preference:
\begin{enumerate}
\item strictly symmetric mechanical layout (identical capacitors, identical mounting,
mirror-image busbar);
\item damping: an RC snubber instead of a plain C, with
$R \approx \sqrt{L_{\mathrm{mesh}}/C}$ as a near-critical starting value, verified by
simulation or measurement;
\item an RCD clamp, which is inherently non-resonant in normal operation.
\end{enumerate}
The dissipation of a damped snubber must be checked: for an RC snubber the resistor
processes approximately $C V^2 f_{sw}$ per charge/discharge cycle pair.

\section{Summary}
\begin{itemize}
\item Snubbers are useful when paralleling 5SNA~1000G650300 modules, primarily as
per-module decoupling capacitors that shorten and equalize the commutation loops.
\item The dimensioning transient is normally the diode reverse recovery (and, for
clamps, short-circuit turn-off), not the nominal IGBT turn-off, which is slow by
design.
\item The voltage budget must be evaluated at chip level, and the cosmic-ray derating
means the apparent margin to \SI{6.5}{kV} is not freely spendable on overshoot.
\item Snubbers do not fix static sharing or gate-side interactions: individual gate
resistors, per-module $C_{GE}$, batch-matched devices, common heatsink, and symmetric
layout remain mandatory.
\item Identical snubbers with identical mounting on every module; otherwise the
snubbers themselves become a source of asymmetry and inter-module ringing.
\end{itemize}

\begin{thebibliography}{9}
\bibitem{ds} ABB Power Grids / Hitachi, ``5SNA~1000G650300 HiPak IGBT module,''
Datasheet 5SYA~1465-03, Oct.~2020.
\bibitem{an2098} ABB Switzerland Ltd, Semiconductors, ``Paralleling of IGBT modules,''
Application Note 5SYA~2098.
\bibitem{an2053} ABB Switzerland Ltd, Semiconductors, ``Applying IGBT,''
Application Note 5SYA~2053.
\bibitem{an2042} ABB Switzerland Ltd, Semiconductors, ``Failure rates of IGBT modules
due to cosmic rays,'' Application Note 5SYA~2042.
\bibitem{snubnote} D.~Bagnara, ``Snubber design guidelines for a high-current IGBT
half-bridge leg,'' internal technical note.
\end{thebibliography}

\end{document}
